\chapter{The particles, and the structure the exact numbers keep}

\textbf{The particles here are put in, and the structure is read out.} The confirmed content of the
Standard Model, six quarks, six leptons and five force carriers, comes from the Particle Data Group at
CERN, and nothing in this chapter derives a particle from first principles. What the chapter shows is
how the exact quantum numbers organize themselves once the masses are set aside, and the reading that
does it is the same equality the shells were read with.

\section{Charge in thirds, and no mass}

Electric charge is quantized in thirds of the electron charge. An up quark is $+2$, a down quark is
$-1$, a charged lepton is $-3$, a neutrino is $0$ and a $W$ is $\pm 3$, all in units of a third, and
carried that way each is an exact integer. Spin is doubled in the same manner, a fermion reading $1$
and a gauge boson $2$, and baryon number is carried in thirds. The ledger at
\texttt{src/\allowbreak{}engine/\allowbreak{}python/\allowbreak{}representation/\allowbreak{}particle/\allowbreak{}standard\_model.py}
holds all of it as integers.

No mass enters. A mass is a measured real number, and reading structure from it would mean choosing a
tolerance, a bound on how close two masses must be to count as one. This work refuses a bound
everywhere, and it refuses one here: every question the chapter asks is answered by equality or by an
exact sum, never by a threshold. Two particles carry the same charge or they do not; a generation's
charges sum to zero or they do not.

\section{The generations recur}

Strip a fermion's generation and its mass, and its signature remains: its charge, its spin, its
color count, its baryon and lepton numbers. The up, charm and top quarks share one signature. The three
charged leptons share another. Read over the three generations, the signature sets are equal. The
generation is one structure seen three times, the way a periodic group was one signature seen down a
column in the second chapter. The only difference between an electron and a muon this reading can see is
a mass, and the reading does not look at it.

That recurrence is the periodic table's lesson at a lower level. There the valence signature repeated
down a column and the instrument read the group without being told chemistry. Here the fermion
signature repeats across a generation and the instrument reads the family without being told a mass.

\section{The charge sum vanishes}

Sum the electric charge of one generation, counting each quark once for each of its three colors:
\[
3 \cdot (+2) \;+\; 3 \cdot (-1) \;+\; (-3) \;+\; 0 \;=\; 0,
\]
in thirds, and the same for the second generation and the third. The sum is exactly zero. That is the
anomaly-free condition the theory needs to be consistent, and it is the reason an atom is neutral: the
three quarks of a proton carry $+3$ thirds and the electron carries $-3$, and they balance to the digit.
The thread from the particle ledger to the atom ledger of the earlier chapters is this exact zero. An
atom is neutral because a generation's charges cancel, and they cancel exactly because charge is an
integer in thirds and the color count is exactly three.

\section{Why a quark is never alone}

A charge that divides by three is an integer charge. The leptons and the bosons all carry one, and they
are the particles seen free. The six quarks do not: their charges are $+2$ and $-1$ in thirds, which no
count of three divides. A quark alone would carry a fractional charge, and none has been seen. Only
combinations whose thirds add to a multiple of three run free, the proton at $+3$ thirds and the
neutron at $0$, and the exact arithmetic of the thirds is all of it. Confinement has more to it than
charge, but the charge alone already forbids a lone quark, and it forbids it by divisibility and not by
any bound.

\section{The hadrons the quarks make}

The combinations that run free are read in
\texttt{examples/\allowbreak{}particle\_physics/\allowbreak{}1\_represent/\allowbreak{}hadrons\_from\_quarks.py}.
A baryon is three quarks and a meson is a quark and an antiquark, and over the light quarks u, d and s
the ten three-quark combinations carry the charges of the baryon decuplet, the proton at $+1$, the
neutron at $0$ and the omega at $-1$ among them, while the nine quark-antiquark combinations carry the
pion and kaon charges. Every one of the nineteen carries an integer charge, a multiple of three thirds,
and a single quark never does. The observed hadron charges are the exact integer sums of their quark
charges, put in from the particle data and added up here, and the proton at $+1$ balancing the electron
is again the neutrality of the atom, read now from the quark side. Nothing is derived about why the
quarks bind or what a hadron weighs; the charge is a sum of integers.

\section{The rarest is the last found}

The census ranks the quantum numbers by magnitude, the total less a count, the rarest first, the same
order the sift kernel probes a field in. Over the eighteen particles the spins fall $1/2$ twelve times,
$1$ five times and $0$ once. The scalar is rarest, and a sift over the particles reaches the Higgs
first, the particle the collider confirmed last. The reading that ranks it first knows nothing of when it was
found; it knows only that a spin of zero occurs once.
